\section*{Bab \ref{ch:b3:spectral} --- Operator Kompak
dan Teorema Spektral}

\begin{solution}{exo:b3:spectral:1}
(a) Himpunan $T(B)$ merupakan himpunan bagian terbatas dari
$\operatorname{im}T$ yang berdimensi berhingga: jadi \omterm{def:b3:spectral:compact}{kompak} relatif menurut Heine--Borel
(\cref{cor:b3:topology:heineborel}, yang dipindahkan lewat \omterm{def:b3:topology:continuity}{homeomorfisma}
linear dengan $\R^n$).
(b) Bahwa $I$ \omterm{def:b3:spectral:compact}{kompak} berarti bola satuan tertutupnya \omterm{def:b3:spectral:compact}{kompak}, yang
menurut teorema Riesz (Tahun ke-2) terjadi tepat dalam dimensi
berhingga. Jika sebuah $T$ yang \omterm{def:b3:spectral:compact}{kompak} berbalikan terbatas $T^{-1}$, maka
$I = T^{-1}T$ akan \omterm{def:b3:spectral:compact}{kompak}
(\cref{prop:b3:spectral:ideal}): yang mustahil dalam dimensi
tak berhingga.
\end{solution}

\begin{solution}{exo:b3:spectral:2}
(a) $\norm{Tx}^2 = \sum\abs{\lambda_n}^2\abs{x_n}^2 \leq
\sup\abs{\lambda_n}^2\norm x^2$, dengan kesamaan hampiran pada
$e_n$ yang mewujudkan supremumnya.
(b) ($\Leftarrow$) Pancungannya $T_N$ (yang menyimpan $n \leq N$ dan
nol selebihnya) mempunyai rank berhingga dengan $\vertiii{T - T_N} =
\sup_{n>N}\abs{\lambda_n} \to 0$: jadi \omterm{def:b3:spectral:compact}{kompak} menurut
\cref{prop:b3:spectral:ideal}. ($\Rightarrow$) Jika
$\abs{\lambda_{n_k}} \geq \delta > 0$ sepanjang sebuah subbarisan:
maka $\norm{Te_{n_k} - Te_{n_l}}^2 = \abs{\lambda_{n_k}}^2 +
\abs{\lambda_{n_l}}^2 \geq 2\delta^2$: sehingga tak ada subbarisan
$(Te_{n_k})$ yang konvergen.
(c) $T^* =$ diagonal dengan $(\bar\lambda_n)$: jadi \omterm{def:b3:spectral:selfadjoint}{swa-adjoin} jika dan hanya jika
semua $\lambda_n \in \R$. Lalu basis bakunya $(e_n)$ merupakan
basis ortonormal berisi vektor eigen, bernilai eigen $\lambda_n \to
0$: yakni teorema spektralnya kata demi kata.
\end{solution}

\begin{solution}{exo:b3:spectral:3}
Misalkan $(x_n)$ terbatas. Maka $(Bx_n)$ terbatas
(sebab $\vertiii B < \infty$); lalu kekompakan $S$ memetik
$SBx_{n_k} \to y$; dan \omterm{def:b3:topology:continuity}{kekontinuan} $A$ memberikan $ASBx_{n_k} \to
Ay$: sehingga $ASB$ \omterm{def:b3:spectral:compact}{kompak}. Jika $ST = TS = I$ dengan $T$ \omterm{def:b3:spectral:compact}{kompak} dan
$\dim H = \infty$: maka $I = ST$ akan \omterm{def:b3:spectral:compact}{kompak}, yang bertentangan dengan
\cref{exo:b3:spectral:1}(b) --- yakni pernyataan yang sama
dilihat dari sisi lainnya.
\end{solution}

\begin{solution}{exo:b3:spectral:4}
(a) Lewat Cauchy--Schwarz pada $y$: $\abs{T_kf(x)}^2 \leq
\bigl(\int\abs{k(x,y)}^2\dd y\bigr)\norm f_2^2$; lalu integralkan terhadap
$x$ (lewat Tonelli): $\norm{T_kf}_2 \leq \norm k_{L^2(\square)}
\norm f_2$.

(b) Keluarga $e_{mn}(x,y) = e_m(x)\overline{e_n(y)}$ bersifat
ortonormal di $L^2(\intcc01^2)$ (sebab Tonelli memisahkan integral
rangkapnya). Untuk ketotalannya: jika $h \perp$ semua $e_{mn}$, maka untuk setiap
$m$, fungsi $y \mapsto \int h(x,y)\overline{e_m(x)}\dd
x$ (yang berada di $L^2$ menurut Cauchy--Schwarz dan Tonelli) ortogonal terhadap
setiap $\overline{e_n}$ --- dan konjugatnya
$(\overline{e_n})$ membentuk \omterm{def:b3:hilbert:onb}{basis Hilbert} setiap kali $(e_n)$ demikian
(sebab pengonjugatan merupakan bijeksi isometrik $L^2$ yang mengawetkan
keortogonalan dan ketotalan) --- sehingga ia $0$ hampir di mana-mana; lalu untuk hampir setiap $y$, $h(\cdot, y) \perp$ setiap
$e_m$: jadi $h(\cdot, y) = 0$ hampir di mana-mana, sehingga $h = 0$ (lewat Tonelli). Jadi
$(e_{mn})$ merupakan \omterm{def:b3:hilbert:onb}{basis Hilbert}; lalu uraikan $k = \sum c_{mn}e_{mn}$.
Pancungannya $k_N$ (dengan indeks $\leq N$) memberikan $T_{k_N}$ yang
mempunyai rank berhingga (dengan jangkauan di $\operatorname{Vect}(e_1, \dots,
e_N)$), dan menurut (a),
\[
\vertiii{T_k - T_{k_N}} \leq \norm{k - k_N}_{L^2} \to 0 :
\]
sehingga $T_k$ merupakan limit \omterm{def:b3:banach:operator}{norma operator} rank berhingga: jadi \omterm{def:b3:spectral:compact}{kompak}
(\cref{prop:b3:spectral:ideal}).
\end{solution}

\begin{solution}{exo:b3:spectral:5}
(a) $\lambda\norm x^2 = \langle x, Tx\rangle \geq 0$ pada sebuah
vektor eigen. Sedangkan rumusnya adalah
\cref{prop:b3:spectral:sanorm} dengan semua nilainya $\langle x,
Tx\rangle \geq 0$: sehingga nilai mutlaknya berlebihan.
(b) Pemetaan $(x, y) \mapsto \langle x, Ty\rangle$ merupakan bentuk
seskuilinear Hermite positif (yang mungkin merosot); dan bukti
Cauchy--Schwarz yang lazim (yakni menguraikan $\langle x + t\eu^{\iu\theta}y,
T(x + t\eu^{\iu\theta}y)\rangle \geq 0$ lalu mengambil
diskriminannya) tak pernah memakai kedefinitannya.
\end{solution}

\begin{solution}{exo:b3:spectral:6}
(a) $\norm{Mf}_2 \leq \norm f_2$, dan pada $f_n = \sqrt n\,
\mathbf 1_{\intcc{1 - 1/n}1}$ (yakni vektor satuan), $\norm{Mf_n}_2
\geq 1 - \frac1n$: sehingga $\vertiii M = 1$; dan ia \omterm{def:b3:spectral:selfadjoint}{swa-adjoin} sebab
pengalinya real. Untuk nilai eigennya: $xf(x) = \lambda f(x)$ hampir di mana-mana
memaksa $f = 0$ hampir di mana-mana di luar himpunan nol $\{x = \lambda\}$: jadi $f =
0$ di $L^2$.
(b) Dengan $f_n$ yang sama: $\norm{Mf_n - f_n}_2 \leq \frac1n \to
0$, sedangkan $f_n \rightharpoonup 0$ (sebab untuk $g \in L^2$ yang tetap,
$\abs{\langle g, f_n\rangle} \leq \norm{g\,\mathbf
1_{\intcc{1-1/n}1}}_2 \to 0$ lewat kekonvergenan terdominasi). Jika $Mf_{n_k} \to h$ dalam
norma, maka $f_{n_k} \to h$, yang memaksa $h = 0$ (lewat limit lemahnya) padahal
$\norm h = 1$: jadi tak ada subbarisan yang konvergen.
(c) Pada \cref{lem:b3:spectral:existence}, tepatnya
pemetikan ``$Tx_{n_k} \to y$'' yang memakai kekompakannya; sedangkan untuk $M$
barisan pemberi maksimumnya memusat di dekat $x = 1$ dan
petanya konvergen lemah ke $0$, tak pernah dalam norma: sehingga vektor eigen
di puncak jangkauan numeriknya sekadar gagal ada.
\end{solution}

\begin{solution}{exo:b3:spectral:7}
(a) $V = T_k$ dengan $k(x, y) = \mathbf 1_{y < x} \in
L^2(\intcc01^2)$: jadi \omterm{def:b3:spectral:compact}{kompak} (\cref{exo:b3:spectral:4}).
Adjoinnya adalah operator kernel berkernel $\overline{k(y,
x)} = \mathbf 1_{y > x}$: yakni $V^*f(x) = \int_x^1f$; sehingga $V \neq V^*$
(ujilah pada $f = \mathbf 1$).
(b) Jika $Vf = \lambda f$ dengan $\lambda \ne 0$: maka $Vf$ \omterm{def:b3:topology:continuity}{kontinu}
pada $\intcc01$ (lewat kekonvergenan terdominasi pada $\int_0^x f$), sehingga $f
= \frac1\lambda Vf$ mempunyai wakil yang \omterm{def:b3:topology:continuity}{kontinu}; lalu
$Vf$ bersifat $\mathcal C^1$ (menurut teorema fundamental kalkulus untuk
integran \omterm{def:b3:topology:continuity}{kontinu}), sehingga $f$ bersifat $\mathcal C^1$, dan
$\lambda f' = f$ dengan $f(0) = \frac1\lambda Vf(0) = 0$: jadi $f =
C\eu^{x/\lambda}$ dengan $C = f(0) = 0$.
(c) Jadi $V$ \omterm{def:b3:spectral:compact}{kompak} tanpa nilai eigen sama sekali kecuali mungkin
$0$ (sebab $Vf = 0$ memaksa $f = 0$ hampir di mana-mana lewat menurunkan
integralnya --- jadi bahkan $0$ pun tidak): sehingga mesin spektralnya
sungguh menuntut \omterm{def:b3:spectral:selfadjoint}{keswa-adjoinan}, bukan sekadar kekompakan.
\end{solution}

\begin{solution}{exo:b3:spectral:8}
Tulis $x = \sum_ic_ie_i + z$ dengan $z \in \ker T$, sehingga $\langle x,
Tx\rangle = \sum_i\mu_i\abs{c_i}^2$. \emph{Batas bawahnya}: pada
bola satuan $V_k = \operatorname{Vect}(e_1, \dots,
e_k)$, $\langle x, Tx\rangle = \sum_{i\leq
k}\mu_i\abs{c_i}^2 \geq \mu_k$: sehingga maksimum atas $V$ dari minimumnya
bernilai $\geq \mu_k$. \emph{Batas atasnya}: misalkan $\dim V = k$ dan $W_k
= \overline{\operatorname{Vect}}(e_k, e_{k+1}, \dots) +
\ker T$, yang berkodimensi $k - 1$ (sebab \omterm{thm:b3:hilbert:decomposition}{komplemen ortogonalnya} adalah
$V_{k-1}$); maka $V \cap W_k \neq \{0\}$ (sebab pemetaan linear $V \to
H/W_k \cong V_{k-1}$ yang ranknya $\leq k - 1$ berkernel
tak trivial), dan sebuah satuan $x \in V\cap W_k$ mempunyai $\langle x,
Tx\rangle = \sum_{i \geq k}\mu_i\abs{c_i}^2 \leq \mu_k$: sehingga
minimum atas $V$ bernilai $\leq \mu_k$. Digabungkan: yakni rumus pertamanya;
sedangkan yang kedua dibuktikan secara setangkup (ujilah $W = W_k$; dan untuk
$W$ sembarang yang berkodimensi $k-1$, $W \cap V_k \neq 0$ memberikan
sebuah vektor satuan dengan $\langle x, Tx\rangle \geq \mu_k$).
Untuk kemonotonannya: $\langle x, Tx\rangle \leq \langle x,
Sx\rangle$ secara titik demi titik berpindah lewat $\max\min$-nya.
\end{solution}

\begin{solution}{exo:b3:spectral:9}
Tulis ulang $f - \lambda Tf = g$. Untuk $\lambda = 0$: $f = g$,
yang selalu \omterm{def:b3:groups:derived}{terselesaikan} secara tunggal. Untuk $\lambda \neq 0$: ini adalah
$(T - \frac1\lambda)f = -\frac g\lambda$, sehingga menurut \omterm{thm:b3:spectral:fredholm}{alternatif
Fredholm} (\cref{thm:b3:spectral:fredholm}) dengan
nilai eigen $\lambda_n = \bigl((n + \frac12)\pi\bigr)^{-2}$ milik
$T$ (\cref{ex:b3:spectral:minxy}): \omterm{def:b3:groups:derived}{keterselesaian} tunggal untuk setiap
$g$ jika dan hanya jika $\frac1\lambda \neq \lambda_n$ untuk setiap $n$, yakni
\[
\lambda \neq \Bigl(n + \tfrac12\Bigr)^2\pi^2
\qquad (n = 0, 1, 2, \dots).
\]
Pada $\lambda = (n+\frac12)^2\pi^2$ yang kecualian: solusinya
ada jika dan hanya jika $g \perp \sin\bigl((n{+}\frac12)\pi x\bigr)$, dan
lalu tunggal sampai penambahan kelipatan sinus itu.
\end{solution}

\begin{solution}{exo:b3:spectral:10}
(a) $Sx = \lambda x$: dengan membandingkan koordinatnya, $0 = \lambda
x_1$ dan $x_n = \lambda x_{n+1}$; lalu jika $\lambda \ne 0$ maka
$x_1 = 0$ dan secara induktif $x = 0$; sedangkan jika $\lambda = 0$, $Sx = 0$
memaksa $x = 0$ (sebab $S$ isometrik). Jadi tak ada nilai eigen. Sedangkan $S^*x =
\lambda x$ berbunyi $x_{n+1} = \lambda x_n$: sehingga $x =
x_1(1, \lambda, \lambda^2, \dots)$, yang berada di $\ell^2$ tepat bila
$\abs\lambda < 1$: yakni satu cakram terbuka penuh berisi nilai eigen.
(b) $\norm{Se_n - Se_m} = \norm{e_{n+1} - e_{m+1}} = \sqrt2$:
sehingga peta $(e_n)$ yang terbatas tak bersubbarisan Cauchy;
dan demikian pula $S^*e_{n+1} = e_n$. Jadi keduanya tak \omterm{def:b3:spectral:compact}{kompak}.
(c) Bagi \omterm{def:b3:spectral:selfadjoint}{operator swa-adjoin} \omterm{def:b3:spectral:compact}{kompak}, nilai eigennya
menangkap segalanya (\cref{thm:b3:spectral:spectral}); sedangkan dengan membuang
salah satu hipotesisnya, nilai eigennya dapat lenyap sama sekali ($S$,
Volterra) atau memenuhi sebuah cakram ($S^*$): sehingga objek yang tegar adalah
spektrum $\{\lambda : T - \lambda I \text{ tak terbalikkan}\}$,
yang teorinya termasuk kuliah berikutnya.
\end{solution}

\begin{solution}{exo:b3:spectral:11}
(a) $S$ merupakan operator diagonal berkoefisien
$\sqrt{\mu_n} \to 0$: jadi \omterm{def:b3:spectral:compact}{kompak}
(\cref{exo:b3:spectral:2}(b), yang dipindahkan ke basis
$(e_n)$ yang dilengkapkan oleh $\ker T$, tempat $S = 0$), \omterm{def:b3:spectral:selfadjoint}{swa-adjoin}
(sebab diagonalnya real), positif (sebab $\langle x, Sx\rangle =
\sum\sqrt{\mu_n}\abs{\langle e_n, x\rangle}^2$), dengan $S^2 =
T$ suku demi suku.

(b) Operator $R$ komutatif dengan $T = R^2$; lalu untuk sebuah vektor eigen $x$ milik
$T$ bernilai eigen $\mu$: $T(Rx) = RTx = \mu Rx$, sehingga
ruang eigen berdimensi berhingga $E_\mu = \ker(T - \mu)$ bersifat
stabil-$R$. Pada $E_\mu$, $R$ bersifat setangkup positif dengan $R^2
= \mu\,\mathrm{id}$: sehingga nilai eigennya $\rho$ memenuhi $\rho^2
= \mu$ dengan $\rho \geq 0$: jadi semuanya sama dengan $\sqrt\mu$, dan
operator terdiagonalkan yang bernilai eigen tunggal bersifat skalar:
sehingga $R = \sqrt\mu\,\mathrm{id}$ pada $E_\mu$. Sedangkan pada $\ker T$: $\norm
{Rx}^2 = \langle x, R^2x\rangle = \langle x, Tx\rangle = 0$.
Jadi $R$ bersesuaian dengan $S$ pada $\ker T$ dan pada setiap ruang eigennya,
yang rentang tertutupnya adalah $H$ (menurut teorema spektralnya): sehingga $R = S$.

(c) Operator $\sqrt G$ bekerja sebagai $\frac1{n\pi}$ pada $e_n =
\sqrt2\sin(n\pi x)$: yakni operator kernel dengan
\[
k(x, y) = \sum_{n\geq1}\frac{2\sin(n\pi x)\sin(n\pi
y)}{n\pi},
\]
dengan deretnya konvergen di $L^2(\intcc01^2)$ (sebab koefisiennya
$\frac1{n\pi} \in \ell^2$; dan kernel jumlah parsialnya memberikan
hampiran rank berhingga). Jadi tak diperlukan bentuk tertutup dasar:
sebab sisi spektralnya \emph{adalah} operatornya.
\end{solution}

\begin{solution}{exo:b3:spectral:12}
(a) $T^*T$ bersifat \omterm{def:b3:spectral:compact}{kompak} (sebab hasil kali operator terbatas dengan \omterm{def:b3:spectral:compact}{operator
kompak}, \cref{exo:b3:spectral:3}), \omterm{def:b3:spectral:selfadjoint}{swa-adjoin}
(sebab $(T^*T)^* = T^*T$), dan positif (sebab $\langle x, T^*Tx\rangle =
\norm{Tx}^2$). Untuk kernelnya: $T^*Tx = 0 \Rightarrow \norm{Tx}^2 =
\langle x, T^*Tx\rangle = 0 \Rightarrow Tx = 0$, dan
sebaliknya: sehingga $\ker T^*T = \ker T$. Lalu teorema spektralnya
menyediakan $(e_n)$ yang ortonormal dengan $T^*Te_n = s_n^2e_n$ dan
$s_n > 0$, yang merentang $(\ker T)^\perp$.

(b) $\langle f_m, f_n\rangle = \frac{\langle Te_m,
Te_n\rangle}{s_ms_n} = \frac{\langle e_m,
T^*Te_n\rangle}{s_ms_n} = \frac{s_n^2}{s_ms_n}\delta_{mn} =
\delta_{mn}$. Lalu uraikan $x = x_0 + \sum_n\langle e_n, x\rangle
e_n$ dengan $x_0 \in \ker T$ (lewat \omterm{thm:b3:hilbert:parseval}{Parseval} pada rentang tertutupnya ditambah
kernelnya); lalu dengan menerapkan $T$ yang \omterm{def:b3:topology:continuity}{kontinu}:
\[
Tx = \sum_n\langle e_n, x\rangle\,Te_n =
\sum_ns_n\langle e_n, x\rangle\,f_n,
\]
dengan deretnya konvergen sebab jumlah parsialnya Cauchy
(karena $\norm{\sum_{N<n\leq M}s_n\langle e_n, x\rangle f_n}^2 =
\sum s_n^2\abs{\langle e_n, x\rangle}^2$, yang terdominasi oleh
$\sup_{n>N}s_n^2\cdot\norm x^2$, dan $s_n \to 0$).

(c) $\norm{Tx}^2 = \sum_ns_n^2\abs{\langle e_n, x\rangle}^2
\leq (\max s_n)^2\norm x^2$, yang tercapai pada $e_n$ pemberi
maksimumnya: sehingga $\vertiii T = \max s_n$. Lalu memancung penguraiannya pada rank
$N$ menyisakan operator bernorma $\sup_{n>N}s_n \to 0$: yakni
hampiran rank berhingga. Sedangkan operator Volterra sama sekali tak bernilai
eigen (\cref{exo:b3:spectral:7}), tetapi $V^*V$ demikian:
sebab $V^*Vf(x) = \int_x^1\int_0^tf(s)\,\dd s\,\dd t$ merupakan
operator kernel setangkup positif (berkernel $1 - \max(x,y)$,
yakni kernel Green bertipe dawai), yang pasangan eigennya --- yang dihitung
lewat soal nilai batas $-u'' = \lambda^{-1}u$ dengan $u'(0)
= u(1) = 0$, yakni keluarga \cref{exo:b3:spectral:9} --- memberikan
nilai singular $s_n = \bigl((n + \frac12)\pi\bigr)^{-1}$.
Jadi penguraian nilai singularnya hidup pada \emph{dua} keluarga ortonormal justru
karena $V$ memutar geometri eigennya menjauh: sehingga tanpa vektor eigen,
namun berstruktur diagonal \omterm{prop:b3:galois:perfect}{sempurna} antara dua basis yang berbeda.
\end{solution}

\begin{solution}{pb:b3:spectral:1}
\textbf{1.} \omterm{def:b3:topology:continuity}{Kekontinuannya}: sebab $\min$ dan $\max$ \omterm{def:b3:topology:continuity}{kontinu};
kesetangkupannya: sebab menukar $x, y$ tak menukar $\min(x,y)$ maupun $1 -
\max(x,y)$. Batasnya: $0 \leq g$, dan batas bertipe
$g(x,y) \leq \max\cdot(1-\max)$ memberikan $g \leq \frac14$ (sebab untuk $u
= \max$: $\min \leq u$ sehingga $g \leq u(1-u) \leq \frac14$). Lalu $g
\in L^2(\square)$: jadi Hilbert--Schmidt, sehingga $G$ \omterm{def:b3:spectral:compact}{kompak}
(\cref{exo:b3:spectral:4}); dan kernelnya real setangkup: sehingga $G$
\omterm{def:b3:spectral:selfadjoint}{swa-adjoin}.

\textbf{2.} Dengan memisahkannya di $y = x$:
\[
u(x) = (1 - x)\int_0^x y\,f(y)\,\dd y +
x\int_x^1(1 - y)\,f(y)\,\dd y .
\]
Untuk $f$ yang \omterm{def:b3:topology:continuity}{kontinu}, turunkanlah (lewat aturan hasil kali dan teorema
fundamentalnya):
\[
u'(x) = -\int_0^xyf + \int_x^1(1-y)f
\qquad\text{(sebab suku batasnya saling meniadakan)},
\]
dan $u''(x) = -xf(x) - (1 - x)f(x) = -f(x)$; sedangkan jelaslah $u(0) =
u(1) = 0$. Sebaliknya jika $u \in \mathcal C^2$ lenyap di kedua
ujungnya, maka $w = u - G(-u'')$ memenuhi $w'' = 0$ dengan $w(0) = w(1) =
0$: sehingga $w$ afin dan lenyap dua kali, jadi $w = 0$.

\textbf{3.} Misalkan $Gf = 0$ dengan $f \in L^2$. Untuk $\psi \in \mathcal
C_c^\infty(\intoo01)$: $\psi = G(-\psi'')$ menurut pertanyaan 2, sehingga
\[
\langle f, \psi\rangle = \langle f, G(-\psi'')\rangle
= \langle Gf, -\psi''\rangle = 0
\]
(sebab $G$ \omterm{def:b3:spectral:selfadjoint}{swa-adjoin}): sehingga menurut lema fundamentalnya
(\cref{cor:b3:lp:fundlemma}), $f = 0$ hampir di mana-mana. Untuk kepositifannya: bagi
$f$ yang \omterm{def:b3:topology:continuity}{kontinu}, dengan $u = Gf$,
\[
\langle f, Gf\rangle = \int_0^1 fu = \int_0^1(-u'')u
= \bigl[-u'u\bigr]_0^1 + \int_0^1(u')^2 = \int_0^1(u')^2
\geq 0 ;
\]
sedangkan untuk $f \in L^2$, hampirilah di $L^2$ oleh $f_n$ yang \omterm{def:b3:topology:continuity}{kontinu}:
maka kedua ruasnya beralih ke limitnya (sebab $G$ terbatas).

\textbf{4.} Jika $Gu = \lambda u$ dengan $\lambda \neq 0$: maka $Gu$
\omterm{def:b3:topology:continuity}{kontinu} (sebab $\abs{Gu(x) - Gu(x')} \leq \norm{g(x,\cdot) -
g(x',\cdot)}_2\norm u_2$, dan kernelnya \omterm{def:b3:topology:continuity}{kontinu} seragam),
sehingga $u$ mempunyai wakil yang \omterm{def:b3:topology:continuity}{kontinu}; lalu
rumus pertanyaan 2 menunjukkan $Gu \in \mathcal C^2$, sehingga $u =
\frac1\lambda Gu \in \mathcal C^2$ dengan $-\lambda u'' =
-(Gu)'' = u$ dan $u(0) = u(1) = 0$.

\textbf{5.} Dari $-\lambda u'' = u$ dengan $u(0) = 0$: $u = A\sin(x/\sqrt
\lambda)$ (dengan $\lambda$ yang positif: sebab menurut pertanyaan 3, $\lambda =
\langle u, Gu\rangle/\norm u^2 > 0$ pada vektor eigennya). Lalu $u(1) =
0$ memaksa $\frac1{\sqrt\lambda} = n\pi$: sehingga $\lambda_n =
\frac1{n^2\pi^2}$, dengan fungsi eigen $\sin(n\pi x)$, yang ternormalkan
$e_n = \sqrt2\sin(n\pi x)$ (sebab $\int_0^12\sin^2(n\pi x)\dd x =
1$). Untuk pemeriksaan keortogonalannya: $2\sin(m\pi x)\sin(n\pi x) =
\cos((m-n)\pi x) - \cos((m+n)\pi x)$ berintegral $0$ untuk $m
\neq n$.

\textbf{6.} Karena $\ker G = \{0\}$ (pertanyaan 3), maka
\cref{thm:b3:spectral:spectral}(1) memberikan $H =
\overline{\operatorname{Vect}}(e_n)$: sehingga sinusnya merupakan \omterm{def:b3:hilbert:onb}{basis
Hilbert} $L^2(\intcc01)$. Untuk $f(x) = x(1 - x)$:
\[
c_n = \sqrt2\int_0^1x(1-x)\sin(n\pi x)\,\dd x
= \sqrt2\;\frac{2\bigl(1 - (-1)^n\bigr)}{n^3\pi^3}
= \begin{cases}\dfrac{4\sqrt2}{n^3\pi^3} & n \text{ gasal},\\
0 & n \text{ genap},\end{cases}
\]
(lewat dua kali pengintegralan parsial). Lalu \omterm{thm:b3:hilbert:parseval}{Parseval}: $\int_0^1x^2(1-x)^2\dd
x = \frac1{30} = \sum_{n \text{ gasal}}\frac{32}{n^6\pi^6}$,
yakni $\sum_{n\text{ gasal}}n^{-6} = \frac{\pi^6}{960}$.

\textbf{7.} Berlaku $\langle e_n, Ge_n\rangle = \lambda_n\norm{e_n}^2
= \lambda_n$. Untuk $x$ yang tetap, koefisien $g(x,
\cdot)$: $\langle e_n, g(x,\cdot)\rangle = (Ge_n)(x) =
\lambda_ne_n(x)$, sehingga $g(x,\cdot) =
\sum_n\lambda_ne_n(x)\,e_n$ di $L^2$. Sedangkan deret eksplisitnya
$\sum_n\lambda_ne_n(x)e_n(y) = \sum_n\frac{2\sin(n\pi
x)\sin(n\pi y)}{n^2\pi^2}$ konvergen secara normal pada bujur sangkarnya
(sebab $\abs{\text{sukunya}} \leq \frac2{n^2\pi^2}$): sehingga jumlahnya
\omterm{def:b3:topology:continuity}{kontinu}, dan untuk setiap $x$ ia berkoefisien
$L^2(\dd y)$ sama dengan $g(x, \cdot)$: jadi kedua fungsi \omterm{def:b3:topology:continuity}{kontinunya}
bersesuaian untuk setiap $(x, y)$. Lalu dengan menetapkan $y = x$ dan
mengintegralkannya (sebab kekonvergenan normalnya mengizinkan pengintegralan yang suku demi suku):
\[
\int_0^1g(x,x)\,\dd x =
\sum_n\lambda_n\int_0^1e_n(x)^2\dd x = \sum_n\lambda_n .
\]

\textbf{8.} Berlaku $\int_0^1g(x,x)\dd x = \int_0^1x(1 - x)\dd x =
\frac16$, sehingga $\sum_{n\geq1}\frac1{n^2\pi^2} = \frac16$:
\[
\zeta(2) = \sum_{n\geq1}\frac1{n^2} = \frac{\pi^2}6 .
\]

\textbf{9.} Untuk $f = \mathbf 1$: $c_n =
\sqrt2\int_0^1\sin(n\pi x)\dd x = \sqrt2\,\frac{1 -
(-1)^n}{n\pi}$: sehingga $c_n = \frac{2\sqrt2}{n\pi}$ untuk $n$ yang gasal, dan $0$
untuk yang genap. Lalu \omterm{thm:b3:hilbert:parseval}{Parseval}: $1 = \sum_{n\text{
gasal}}\frac{8}{n^2\pi^2}$, sehingga $\sum_{n\text{ gasal}}n^{-2} =
\frac{\pi^2}8$, dan $\zeta(2) = \frac{\pi^2}8\cdot\frac{1}{1
- \frac14} = \frac{\pi^2}6$ (sebab suku genapnya adalah $\frac14\zeta(2)$).
Perbandingannya: \omterm{thm:b3:hilbert:parseval}{Parseval} untuk satu $f$ menjumlahkan
$\abs{\langle e_n, f\rangle}^2$; sedangkan rumus tracenya mengintegralkan
diagonal \emph{kernelnya}, yang setara dengan menjumlahkan
\omterm{thm:b3:hilbert:parseval}{Parseval} atas seluruh keluarga ortonormalnya sekaligus ---
$\sum_n\langle e_n, Ge_n\rangle$ --- sehingga ia buta terhadap
pilihan fungsi uji tertentu mana pun.

\textbf{10.} Berlaku $\abs{c_n\cos(n\pi t)} \leq \abs{c_n}$ dengan
$\sum\abs{c_n}^2 < \infty$: sehingga untuk setiap $t$ deretnya konvergen
di $L^2$ (lewat uraian ortonormalnya,
\cref{thm:b3:hilbert:parseval}(3)); sedangkan batas ekornya
$\norm{u(t) - u_N(t)}_2^2 \leq \sum_{n>N}\abs{c_n}^2$ bersifat
seragam terhadap $t$, dan setiap jumlah parsialnya \omterm{def:b3:topology:continuity}{kontinu} terhadap $t$
(sebab berupa berhingga banyak kosinus): jadi $t \mapsto u(t,\cdot)$ \omterm{def:b3:topology:continuity}{kontinu}
ke dalam $L^2$. Untuk $f = \sum_{n\leq N}c_ne_n$: setiap ragamnya
$\cos(n\pi t)\sin(n\pi x)$ memenuhi $\partial_t^2 =
-n^2\pi^2 = \partial_x^2$ yang diterapkan padanya, lenyap di $x = 0,
1$, bernilai $\sin(n\pi x)$ dan berturunan waktu $0$ di $t =
0$: sehingga jumlah berhingganya menyelesaikan segalanya. Secara musik: gerak
dawainya merupakan tumpangan \emph{gelombang tegak} $e_n$,
yang frekuensinya $n\pi$ adalah nada dasar dan
nada atasnya; sedangkan teorema spektralnya mengatakan bahwa setiap bentuk awal
terurai secara tunggal menjadi nada murni ini, dengan koefisien
$c_n$ sebagai warna suaranya. Jadi mendengar dawai adalah menghitung
uraian ortonormal.

\textbf{11.} Dengan $a_n = \abs{\langle u_n, x\rangle}^2$:
$m_{p+1} = \sum_n\mu_n^{p+1}a_n = \sum_n\bigl(\mu_n^{p/2}
\sqrt{a_n}\bigr)\bigl(\mu_n^{p/2+1}\sqrt{a_n}\bigr) \leq
\sqrt{m_p\,m_{p+2}}$ (lewat Cauchy--Schwarz di $\ell^2$). Karena itu
nisbah $m_{p+1}/m_p$ bersifat tak turun terhadap $p$; lalu karena
$R(x) = \frac{m_1}{m_0}$, $\frac{\langle Ax, Ax\rangle}
{\langle x, Ax\rangle} = \frac{m_2}{m_1}$ dan $R(Ax) =
\frac{m_3}{m_2}$, rantainya menyusul, dengan setiap sukunya $\leq \mu_1$
sebab $m_{p+1} \leq \mu_1m_p$ suku demi suku. Untuk kekonvergenannya: jika
$a_1 > 0$ (dengan menulis bobot total nilai eigen puncaknya sebagai
$a_1$), maka
\[
\mu_1 \geq R(A^kx) = \frac{m_{2k+1}}{m_{2k}} =
\mu_1\,\frac{a_1 + \sum_{\mu_n<\mu_1}(\mu_n/\mu_1)^{2k+1}
a_n}{a_1 + \sum_{\mu_n<\mu_1}(\mu_n/\mu_1)^{2k}a_n}
\longrightarrow \mu_1,
\]
lewat kekonvergenan terdominasi jumlahnya (sebab nisbahnya $< 1$): sehingga
metode pangkatnya konvergen untuk setiap vektor awal yang tak
ortogonal terhadap ruang eigen puncaknya.

\textbf{12.} Secara titik demi titik, $\abs{\langle x, (A - B)x\rangle}
\leq \vertiii{A - B}\,\norm x^2$, sehingga $R_A(x) \leq R_B(x) +
\vertiii{A - B}$ untuk setiap $x$. Lalu memasukkan ini ke
rumus maks--min \cref{exo:b3:spectral:8}: $\mu_n(A)
\leq \mu_n(B) + \vertiii{A - B}$, dan secara setangkup pada $A,
B$: sehingga $\abs{\mu_n(A) - \mu_n(B)} \leq \vertiii{A - B}$ untuk
setiap $n$ sekaligus.

\textbf{13.} Persamaan $-w'' = x - x^2$ \omterm{def:b3:lebesgue:l1}{terintegralkan} menjadi $w = -\frac{x^3}6
+ \frac{x^4}{12} + cx$ (dengan $w(0) = 0$), dan $w(1) = 0$
memberikan $c = \frac1{12}$:
\[
w = \frac{x^4 - 2x^3 + x}{12} = \frac{x(1-x)(1 + x -
x^2)}{12} = Gu .
\]
Lalu $\norm u_2^2 = \int_0^1x^2(1-x)^2 = \frac1{30}$, dan
dengan $\int_0^1x^3(1-x)^3 = B(4,4) = \frac1{140}$:
\[
\langle u, Gu\rangle = \frac1{12}\Bigl(\frac1{30} +
\frac1{140}\Bigr) = \frac{17}{5040},
\qquad
R(u) = \frac{17/5040}{1/30} = \frac{17}{168} .
\]
Jadi $\frac1{\pi^2} = \lambda_1 \geq \frac{17}{168}$, yakni
$\pi^2 \leq \frac{168}{17} = 9.8824$: sehingga $\pi \leq 3.14364 <
3.1437$ (dengan nilai sejatinya $\pi^2 = 9.8696$). Yakni sebuah polinomial, sebuah
integral, sebuah angka.

\textbf{14.} Tulis $u = x - x^2$, sehingga $Gu =
\frac{u(1 + u)}{12}$ dan, dengan memakai $\int u^2 = \frac1{30}$,
$\int u^3 = \frac1{140}$, dan $\int u^4 = B(5,5) =
\frac{4!\,4!}{9!} = \frac1{630}$:
\[
\norm{Gu}_2^2 = \frac1{144}\int u^2(1+u)^2 =
\frac1{144}\Bigl(\frac1{30} + \frac2{140} +
\frac1{630}\Bigr) = \frac1{144}\cdot\frac{62}{1260} =
\frac{31}{90720} .
\]
Karena itu $\frac{\langle Gu, Gu\rangle}{\langle u, Gu\rangle} =
\frac{31/90720}{17/5040} = \frac{31}{306}$, dan menurut pertanyaan
11 ini tetap $\leq \mu_1 = \frac1{\pi^2}$: sehingga $\pi^2 \leq
\frac{306}{31} = 9.87097$, yakni $\pi \leq 3.14181 <
3.1419$ --- yakni empat angka (dan iterasi berikutnya akan memberikan
sekitar delapan, sebab galatnya mengerut sebesar
$(\lambda_2/\lambda_1)^2 = \frac1{16}$ setiap langkah).

\textbf{15.} Keluarga $(e_m \otimes e_n)(x,y) =
e_m(x)e_n(y)$ merupakan \omterm{def:b3:hilbert:onb}{basis Hilbert} $L^2(\intcc01^2)$
(dengan keortonormalannya lewat Tonelli; dan ketotalannya seperti pada
\cref{exo:b3:spectral:5}). Menurut pertanyaan 7, untuk $x$ yang tetap:
$g(x, \cdot) = \sum_n\lambda_ne_n(x)e_n$, sehingga koefisien
$g$ pada $e_m\otimes e_n$ adalah
\[
\langle e_m\otimes e_n,\ g\rangle
= \int_0^1 e_m(x)\,\lambda_n e_n(x)\,\dd x
= \lambda_n\,\delta_{mn} .
\]
Lalu \omterm{thm:b3:hilbert:parseval}{Parseval} pada bujur sangkarnya:
\[
\iint g^2 = \sum_{m,n}\abs{\langle e_m\otimes e_n,
g\rangle}^2 = \sum_n\lambda_n^2 .
\]

\textbf{16.} Menurut kesetangkupan $g$,
\[
\iint g^2 = 2\iint_{y<x}y^2(1-x)^2
= 2\int_0^1(1-x)^2\,\frac{x^3}3\,\dd x
= \frac23\,B(4, 3) = \frac23\cdot\frac{3!\,2!}{6!}
= \frac23\cdot\frac1{60} = \frac1{90} .
\]

\textbf{17.} Digabungkan: $\sum_n\frac1{n^4\pi^4} =
\frac1{90}$, yakni $\zeta(4) = \frac{\pi^4}{90}$. Secara
umum, $\operatorname{tr}(G^k) = \sum\lambda_n^k =
\frac{\zeta(2k)}{\pi^{2k}}$ sama dengan sebuah integral teriterasi atas
hasil kali kernel polinomial rasional $g$ pada
kubus-$k$: yakni sebuah bilangan rasional. Karena itu $\zeta(2k) \in
\pi^{2k}\Q$ untuk setiap $k$. Mesinnya hanya mencapai argumen
genap sebab nilai eigennya masuk lewat \emph{pangkatnya}
--- yakni $\sum\lambda_n^k$ --- dan $\lambda_n =
\frac1{n^2\pi^2}$: sehingga tak ada gabungan trace yang menghasilkan $\sum
n^{-3}$; sedangkan sifat aritmetika $\zeta(3)$ (yang irasional menurut
Apéry, dengan ketranssendenannya masih terbuka) terletak di luar
pembukuan spektralnya.

\textbf{18.} Suku puncak sebuah jumlah bersuku positif bernilai paling
banyak jumlahnya: $\lambda_1^2 \leq \sum\lambda_n^2 =
\frac1{90}$, sehingga $\frac1{\pi^2} \leq \frac1{\sqrt{90}}$ dan
$\pi \geq 90^{1/4} = 3.0801\ldots$ Bersama pertanyaan 14:
$3.080 < \pi < 3.1419$, lewat aritmetika dawai belaka. Sedangkan trace
yang lebih tinggi menajamkan batas bawahnya secara geometri:
$\lambda_1 \leq (\operatorname{tr}G^{2k})^{1/2k} =
\lambda_1\bigl(1 + \sum_{n\geq2}(\lambda_n/\lambda_1)^{2k}
\bigr)^{1/2k}$, dan faktor parasitnya mati secepat
$\bigl(\tfrac14\bigr)^{2k}\cdot\frac1{2k}$ --- yakni mekanisme
jurang spektral yang sama seperti kekonvergenan metode pangkat
(pertanyaan 11), yang dilihat dari sisi tracenya.

\textbf{19.} Berlaku $\abs{n^2\pi^2 - \nu} \geq \delta > 0$ untuk setiap
$n$ (sebab barisan $n^2\pi^2 \to \infty$ menghindari $\nu$ dengan sebuah
selisih), dan $\abs{n^2\pi^2 - \nu} \geq \frac{n^2\pi^2}2$
untuk $n$ yang besar. Untuk kekonvergenan $L^2$-nya: koefisiennya
$\frac{c_n}{n^2\pi^2 - \nu}$ terjumlahkan kuadrat (sebab terdominasi
oleh $\frac{\abs{c_n}}\delta$). Untuk kekonvergenan seragamnya: norma sup
ekornya terbatas oleh $\sqrt2\sum_{n>N}
\frac{\abs{c_n}}{\abs{n^2\pi^2 - \nu}} \leq
\frac{2\sqrt2}{\pi^2}\bigl(\sum\abs{c_n}^2\bigr)^{1/2}
\bigl(\sum_{n>N}n^{-4}\bigr)^{1/2} \to 0$ (lewat Cauchy--Schwarz).
Untuk pemeriksaannya: $Gf + \nu Gu$ berkoefisien-$e_n$
\[
\lambda_nc_n + \frac{\nu\lambda_nc_n}{n^2\pi^2 - \nu}
= \frac{c_n}{n^2\pi^2}\Bigl(1 + \frac{\nu}{n^2\pi^2 -
\nu}\Bigr) = \frac{c_n}{n^2\pi^2 - \nu} :
\]
yakni tepat
koefisien $u$, sehingga $u = Gf + \nu Gu$; sedangkan ketunggalannya
sebab selisih $v$ dua solusinya memenuhi $v = \nu
Gv$, yakni $\langle e_n, v\rangle(n^2\pi^2 - \nu) = 0$ untuk
setiap $n$: jadi $v = 0$.

\textbf{20.} Seperti pada pertanyaan 19, persamaan $u = Gf + \nu
Gu$ setara dengan keluarga persamaan koefisiennya
$(n^2\pi^2 - \nu)\,\langle e_n, u\rangle = c_n$ dengan $n \geq
1$. Untuk $n = m$ ruas kirinya adalah
$0$: sehingga \omterm{def:b3:groups:derived}{keterselesaiannya} memaksa $c_m = 0$, dan lalu $\langle e_m,
u\rangle$ bebas sedangkan semua koefisien lainnya
tertentukan: jadi solusinya membentuk garis $u_0 + \R e_m$.
Untuk resonansinya: sebuah penggerak yang berkomponen pada ragam eigennya
memompakan energi ke dalamnya tanpa batas --- yakni ayunan yang didorong pada
frekuensinya sendiri.

\textbf{21.} Dari rumus pertanyaan 19, $\norm{R_\nu f}_2^2
= \sum\frac{\abs{c_n}^2}{(n^2\pi^2 - \nu)^2} \leq
\frac{\norm f^2}{(\pi^2 - \nu)^2}$ (sebab untuk $\nu < \pi^2$
nilai eigen terdekatnya adalah $\pi^2$), dengan kesamaannya didekati pada
$f = e_1$: sehingga \omterm{def:b3:banach:operator}{norma operatornya} $\frac1{\pi^2 - \nu}$. Untuk kekompakannya:
$R_\nu$ merupakan limit norma pancungan rank berhingganya
(sebab koefisien ekornya $\frac1{n^2\pi^2 - \nu} \to 0$);
sedangkan \omterm{def:b3:spectral:selfadjoint}{keswa-adjoinan} dan kepositifannya terbaca dari bentuk
diagonalnya (sebab semua koefisiennya $\frac1{n^2\pi^2 - \nu} > 0$).
Jadi $R_\nu$ bernilai eigen $\frac1{n^2\pi^2 - \nu}$: sehingga
analisis Bagian I--VI berulang kata demi kata.

\textbf{22.} Kamusnya: \emph{nilai eigen} $\lambda_n =
\frac1{n^2\pi^2}$ $\leftrightarrow$ kuadrat frekuensi
$n^2\pi^2$ harmonik ke-$n$; \emph{trace} $\sum
\lambda_n = \frac16$ $\leftrightarrow$ $\zeta(2) =
\frac{\pi^2}6$; \emph{norma Hilbert--Schmidt} $\iint g^2 =
\frac1{90}$ $\leftrightarrow$ $\zeta(4) = \frac{\pi^4}{90}$;
\emph{\omterm{thm:b3:spectral:fredholm}{alternatif Fredholm}} $\leftrightarrow$ resonansi
dawai yang digerakkan; \emph{min--maks} $\leftrightarrow$
taksiran variasi, sampai $\pi < 3.1437$ dari satu
polinomial. Di balik setiap pasangannya, objek yang sama: yakni satu
\omterm{def:b3:spectral:selfadjoint}{operator swa-adjoin} \omterm{def:b3:spectral:compact}{kompak}, yang didiagonalkan sekali lalu dieksploitasi
lima cara.

\textbf{23.} Dengan memisahkannya menurut paritas lalu menyubstitusikan $n = 2m$
pada bagian genapnya,
\[
\zeta(6) = \sum_{n\text{ gasal}}\frac1{n^6} +
\sum_{m\geq1}\frac1{(2m)^6}
= \frac{\pi^6}{960} + \frac{\zeta(6)}{64},
\]
sehingga $\frac{63}{64}\zeta(6) = \frac{\pi^6}{960}$ dan $\zeta(6)
= \frac{64\,\pi^6}{63\cdot960} = \frac{\pi^6}{945}$. Sedangkan
jalur tracenya akan menghitung $\operatorname{tr}G^3 =
\sum\lambda_n^3 = \zeta(6)/\pi^6$ sebagai $\iint g\,g_2$ dengan
kernel teriterasi $g_2(x,y) = \int_0^1g(x,z)g(z,y)\dd z$ ---
yakni tiga pengintegralan polinomial sepenggal-sepenggal; padahal \omterm{thm:b3:hilbert:parseval}{Parseval} pada
$x(1-x)$ hanya memerlukan satu.

\textbf{24.} (a) Diagonalkanlah: $v = \sum_na_nu_n$ (ditambah sebuah
komponen kernel yang mungkin, yang $\langle v, Av\rangle$-nya tak memperoleh
apa pun sedangkan $\norm v^2$-nya tumbuh, sehingga pemberi maksimumnya tak
memilikinya). Maka $\langle v, Av\rangle = \sum\mu_na_n^2 \leq
\mu_1\sum a_n^2$, dengan kesamaannya jika dan hanya jika $a_n = 0$ setiap kali
$\mu_n < \mu_1$: sehingga pemberi maksimumnya terletak di
ruang eigen-$\mu_1$. (b) Untuk sembarang $u$,
\[
\langle\abs u, A\abs u\rangle - \langle u, Au\rangle
= \iint k(x,y)\,\bigl(\abs{u(x)}\abs{u(y)} -
u(x)u(y)\bigr)\dd x\,\dd y \;\geq\; 0,
\]
sebab integrannya tak negatif di setiap titiknya. Jika $P = \{u >
0\}$ dan $N = \{u < 0\}$ keduanya berukuran positif, maka pada
$P\times N$ integrannya sama dengan $2k\abs{u(x)}\abs{u(y)} >
0$ pada himpunan berukuran positif: sehingga ketaksamaannya sejati. Sedangkan
sebuah fungsi eigen-$\mu_1$ $u$ memaksimumkan hasil bagi Rayleighnya,
dan $\abs u$ bernorma sama, sehingga kesejatiannya akan menampilkan
$R(\abs u) > \mu_1$ --- yang mustahil; jadi $u$ bertanda tetap
hampir di mana-mana, katakanlah $u \geq 0$. Lalu $u(x) = \mu_1^{-1}(Au)(x) =
\mu_1^{-1}\int k(x,y)u(y)\dd y > 0$ untuk setiap $x \in
\intoo01$ (sebab $k(x,\cdot) > 0$ dan $u \neq 0$). (c) Jika
ruang eigennya berdimensi $\geq 2$, ia akan memuat dua
fungsi eigen yang \emph{ortogonal} $u, v$, yang masing-masing bertanda tetap
dan tak lenyap di \omterm{def:b3:topology:interior}{interiornya} menurut (b); tetapi lalu
$\abs{\langle u, v\rangle} = \int\abs u\,\abs v > 0$ ---
yang bertentangan. Pada dawainya: $k = g > 0$ pada bujur sangkar
terbukanya, $\mu_1 = \lambda_1 = \frac1{\pi^2}$ memang
sederhana, dan $e_1 = \sqrt2\sin(\pi x)$ positif pada
$\intoo01$; sedangkan setiap $e_n = \sqrt2\sin(n\pi x)$ dengan $n \geq 2$,
yang ortogonal terhadap $e_1$ yang positif, haruslah berintegral nol
terhadapnya, sehingga berganti tanda --- seperti dibenarkan $n - 1$ akar
\omterm{def:b3:topology:interior}{interiornya} $\frac kn$.

\textbf{25.} Karena $n^2\pi^2 \to \infty$, minimumnya $d =
\min_n\abs{n^2\pi^2 - \nu}$ tercapai, pada suatu ragam $m$,
dan $d > 0$ sebab $\nu$ menghindari spektrumnya. Lalu rumus diagonal
pertanyaan 19 memberikan
\[
\norm{R_\nu f}_2^2 =
\sum_n\frac{\abs{c_n}^2}{(n^2\pi^2 - \nu)^2}
\leq \frac1{d^2}\,\norm f_2^2,
\]
dengan kesamaannya untuk $f = e_m$: sehingga $\vertiii{R_\nu} = \frac1d$,
yakni kebalikan jarak dari $\nu$ ke spektrumnya
--- yaitu asas resolven umumnya, di sini dalam koordinat yang
eksplisit. \omterm{def:b3:spectral:selfadjoint}{Keswa-adjoinannya} terbaca dari koefisien diagonalnya
yang real; sedangkan kekompakannya menyusul seperti pada pertanyaan 21 (sebab
koefisiennya menuju $0$, sehingga pancungan rank berhingganya
konvergen dalam norma). Untuk harga resonansinya: bagi $f = e_1$ dan $\nu =
(1-\varepsilon)\pi^2$, rumusnya memberikan $u =
\frac{c_1}{\pi^2 - \nu}e_1 = \frac1{\varepsilon\pi^2}e_1$,
terhadap tanggapan statisnya $Ge_1 = \frac1{\pi^2}e_1$: yakni
penguatan sebesar $\frac1\varepsilon$. Pada $\varepsilon =
10^{-2}$ tanggapannya $100$ kali tanggapan statisnya --- dan
ia divergen saat $\varepsilon \to 0$, yang merupakan alternatif
pertanyaan 20 yang dilihat dari sisi terbatasnya: sehingga makin dekat
frekuensi penggeraknya ke frekuensi alaminya, makin tak terbatas
balikannya.

\end{solution}
